\part{Modern transformer phases}

\chapter{LLaMA-style architectural changes}
\label{ch:llama-style}
After GPT-2 correctness, the roadmap adds LLaMA-style changes such as RoPE, RMSNorm, and gated MLPs. \citep{touvron2023llama}



\begin{intuitionbox}{Intuition}
Modernisation changes contracts, not just modules.
\end{intuitionbox}

\begin{definitionbox}{Mathematical or contract statement}
A modern dense decoder remains autoregressive but may change positions, norms, MLPs, and tensor names.
\end{definitionbox}

\begin{implementationbox}{Implementation interpretation}
Extend config, converter, validator, and tests before adding kernels.
\end{implementationbox}

\begin{verificationbox}{How to verify this works}
\begin{itemize}
\item Check shapes and dtypes at the boundary.
\item Compare deterministic outputs against PyTorch golden vectors.
\item Record tolerance, seed, model artifact, and backend.
\item Do not benchmark until validation passes.
\end{itemize}
\end{verificationbox}

\begin{warningbox}{Common failure modes}
\begin{itemize}
\item Letting framework defaults choose dtype or layout.
\item Testing only a batch size of one.
\item Depending on upstream checkpoint names in runtime code.
\item Treating benchmark output as validation.
\end{itemize}
\end{warningbox}

\begin{chapterexercises}
\item State the central invariant in one sentence.\\
\textit{Hint.} Use the conceptual guide vocabulary.
\item Construct a toy example with small shapes.\\
\textit{Hint.} Use $B=2$, $T=3$, $H=8$ where possible.
\item Name a validation test that would catch a plausible bug.\\
\textit{Hint.} Prefer generated golden vectors to handwritten long arrays.
\end{chapterexercises}

\chapter{RoPE}
\label{ch:rope}
RoPE rotates query and key channel pairs by position and is not enabled for phase-0 GPT-2. \citep{su2021roformer}

\[
\theta_{p,j}=p\,\omega_j,\qquad \omega_j=\theta_{\mathrm{base}}^{-2j/d_h},
\]
and each even/odd channel pair is rotated as
\[
\begin{pmatrix}x'_{2j}\\x'_{2j+1}\end{pmatrix}
=
\begin{pmatrix}\cos\theta_{p,j}&-\sin\theta_{p,j}\\
\sin\theta_{p,j}&\cos\theta_{p,j}\end{pmatrix}
\begin{pmatrix}x_{2j}\\x_{2j+1}\end{pmatrix}.
\]

\begin{intuitionbox}{Intuition}
RoPE moves position information into attention geometry.
\end{intuitionbox}

\begin{definitionbox}{Mathematical or contract statement}
RoPE applies a block-diagonal rotation to Q and K, not V. Decode must use absolute position ids consistent with prefill so cached rotated keys match full recomputation.
\end{definitionbox}

\begin{implementationbox}{Implementation interpretation}
Store cached keys after the same rotation used in full recomputation.
\end{implementationbox}

\begin{verificationbox}{How to verify this works}
\begin{itemize}
\item Check shapes and dtypes at the boundary.
\item Compare deterministic outputs against PyTorch golden vectors.
\item Record tolerance, seed, model artifact, and backend.
\item Do not benchmark until validation passes.
\end{itemize}
\end{verificationbox}

\begin{warningbox}{Common failure modes}
\begin{itemize}
\item Letting framework defaults choose dtype or layout.
\item Testing only a batch size of one.
\item Depending on upstream checkpoint names in runtime code.
\item Treating benchmark output as validation.
\end{itemize}
\end{warningbox}

\begin{chapterexercises}
\item State the central invariant in one sentence.\\
\textit{Hint.} Use the conceptual guide vocabulary.
\item Construct a toy example with small shapes.\\
\textit{Hint.} Use $B=2$, $T=3$, $H=8$ where possible.
\item Name a validation test that would catch a plausible bug.\\
\textit{Hint.} Prefer generated golden vectors to handwritten long arrays.
\end{chapterexercises}

\chapter{RMSNorm}
\label{ch:rmsnorm}
RMSNorm is common in newer decoder families and normalises by root mean square. \citep{zhang2019rmsnorm}

\[
\operatorname{RMSNorm}(x)_j=\gamma_j\frac{x_j}{\sqrt{H^{-1}\sum_{k=1}^{H}x_k^2+\epsilon}}.
\]

\begin{intuitionbox}{Intuition}
RMSNorm scales without mean subtraction.
\end{intuitionbox}

\begin{definitionbox}{Mathematical or contract statement}
RMSNorm divides by the square root of the mean square and multiplies by a learned scale.
\end{definitionbox}

\begin{implementationbox}{Implementation interpretation}
Implement it as a separate module from LayerNorm.
\end{implementationbox}

\begin{verificationbox}{How to verify this works}
\begin{itemize}
\item Check shapes and dtypes at the boundary.
\item Compare deterministic outputs against PyTorch golden vectors.
\item Record tolerance, seed, model artifact, and backend.
\item Do not benchmark until validation passes.
\end{itemize}
\end{verificationbox}

\begin{warningbox}{Common failure modes}
\begin{itemize}
\item Letting framework defaults choose dtype or layout.
\item Testing only a batch size of one.
\item Depending on upstream checkpoint names in runtime code.
\item Treating benchmark output as validation.
\end{itemize}
\end{warningbox}

\begin{chapterexercises}
\item State the central invariant in one sentence.\\
\textit{Hint.} Use the conceptual guide vocabulary.
\item Construct a toy example with small shapes.\\
\textit{Hint.} Use $B=2$, $T=3$, $H=8$ where possible.
\item Name a validation test that would catch a plausible bug.\\
\textit{Hint.} Prefer generated golden vectors to handwritten long arrays.
\end{chapterexercises}

\chapter{Gated MLPs: SwiGLU and GeGLU}
\label{ch:gated-mlps}
Gated MLPs use a learned gate and are common in modern decoder families. \citep{shazeer2020glu}

\[
\operatorname{SwiGLU}(x)=
\left(\operatorname{SiLU}(xW_g^\top+b_g)\odot (xW_u^\top+b_u)\right)W_d^\top+b_d,
\]
where \(W_g,W_u\in\R^{F\times H}\), \(W_d\in\R^{H\times F}\), and biases are present only for model families whose config enables them.

\begin{intuitionbox}{Intuition}
The gate controls which expanded channels flow.
\end{intuitionbox}

\begin{definitionbox}{Mathematical or contract statement}
SwiGLU and GeGLU compute two expanded vectors in \(\R^F\), combine them with an elementwise product, and then apply the down projection back to \(\R^H\).
\end{definitionbox}

\begin{implementationbox}{Implementation interpretation}
Map gate, up, and down tensors separately in the converter.
\end{implementationbox}

\begin{verificationbox}{How to verify this works}
\begin{itemize}
\item Check shapes and dtypes at the boundary.
\item Compare deterministic outputs against PyTorch golden vectors.
\item Record tolerance, seed, model artifact, and backend.
\item Do not benchmark until validation passes.
\end{itemize}
\end{verificationbox}

\begin{warningbox}{Common failure modes}
\begin{itemize}
\item Letting framework defaults choose dtype or layout.
\item Testing only a batch size of one.
\item Depending on upstream checkpoint names in runtime code.
\item Treating benchmark output as validation.
\end{itemize}
\end{warningbox}

\begin{chapterexercises}
\item State the central invariant in one sentence.\\
\textit{Hint.} Use the conceptual guide vocabulary.
\item Construct a toy example with small shapes.\\
\textit{Hint.} Use $B=2$, $T=3$, $H=8$ where possible.
\item Name a validation test that would catch a plausible bug.\\
\textit{Hint.} Prefer generated golden vectors to handwritten long arrays.
\end{chapterexercises}

\chapter{MQA}
\label{ch:mqa}
MQA uses many query heads but one or few key/value heads, reducing decode cache bandwidth. \citep{shazeer2019mqa}



\begin{intuitionbox}{Intuition}
Fewer KV heads mean fewer cache bytes per token.
\end{intuitionbox}

\begin{definitionbox}{Mathematical or contract statement}
Strict MQA is the special case \(H_{kv}=1\) while \(H_q\) may be larger. Implementations can share the same grouped-head machinery as GQA, but the config must still record the exact number of key/value heads.
\end{definitionbox}

\begin{implementationbox}{Implementation interpretation}
Generalise attention code without changing phase-0 GPT-2 equality.
\end{implementationbox}

\begin{verificationbox}{How to verify this works}
\begin{itemize}
\item Check shapes and dtypes at the boundary.
\item Compare deterministic outputs against PyTorch golden vectors.
\item Record tolerance, seed, model artifact, and backend.
\item Do not benchmark until validation passes.
\end{itemize}
\end{verificationbox}

\begin{warningbox}{Common failure modes}
\begin{itemize}
\item Letting framework defaults choose dtype or layout.
\item Testing only a batch size of one.
\item Depending on upstream checkpoint names in runtime code.
\item Treating benchmark output as validation.
\end{itemize}
\end{warningbox}

\begin{chapterexercises}
\item State the central invariant in one sentence.\\
\textit{Hint.} Use the conceptual guide vocabulary.
\item Construct a toy example with small shapes.\\
\textit{Hint.} Use $B=2$, $T=3$, $H=8$ where possible.
\item Name a validation test that would catch a plausible bug.\\
\textit{Hint.} Prefer generated golden vectors to handwritten long arrays.
\end{chapterexercises}

\chapter{GQA}
\label{ch:gqa}
GQA groups query heads over fewer key/value heads. \citep{ainslie2023gqa}

\begin{figure}[htbp]
\centering
\resizebox{0.94\textwidth}{!}{%
\begin{tikzpicture}[>=Latex]
\foreach \i in {0,...,7}{\node[draw,fill=EngineBlue!8,minimum width=7mm,minimum height=7mm] (q\i) at (\i*0.75,0) {$Q_{\i}$};}
\foreach \j in {0,...,1}{\node[draw,fill=EngineTeal!12,minimum width=14mm,minimum height=7mm] (kv\j) at (\j*3.0+1.1,-1.3) {$K,V_{\j}$};}
\foreach \i in {0,1,2,3}{\draw[->,EngineTeal!80!black] (q\i.south)--(kv0.north);}
\foreach \i in {4,5,6,7}{\draw[->,EngineTeal!80!black] (q\i.south)--(kv1.north);}
\end{tikzpicture}
}
\caption{GQA maps groups of query heads to fewer key/value heads.}
\end{figure}


\begin{intuitionbox}{Intuition}
GQA is a compromise between MHA quality and MQA bandwidth.
\end{intuitionbox}

\begin{definitionbox}{Mathematical or contract statement}
If \(g=H_q/H_{kv}\) is an integer, query head \(h\) reads key/value head \(\lfloor h/g\rfloor\). The cache stores \(H_{kv}\) heads, while the attention output still has \(H_q\) query-head outputs.
\end{definitionbox}

\begin{implementationbox}{Implementation interpretation}
Validate grouping with small distinguishable heads.
\end{implementationbox}

\begin{verificationbox}{How to verify this works}
\begin{itemize}
\item Check shapes and dtypes at the boundary.
\item Compare deterministic outputs against PyTorch golden vectors.
\item Record tolerance, seed, model artifact, and backend.
\item Do not benchmark until validation passes.
\end{itemize}
\end{verificationbox}

\begin{warningbox}{Common failure modes}
\begin{itemize}
\item Letting framework defaults choose dtype or layout.
\item Testing only a batch size of one.
\item Depending on upstream checkpoint names in runtime code.
\item Treating benchmark output as validation.
\end{itemize}
\end{warningbox}

\begin{chapterexercises}
\item State the central invariant in one sentence.\\
\textit{Hint.} Use the conceptual guide vocabulary.
\item Construct a toy example with small shapes.\\
\textit{Hint.} Use $B=2$, $T=3$, $H=8$ where possible.
\item Name a validation test that would catch a plausible bug.\\
\textit{Hint.} Prefer generated golden vectors to handwritten long arrays.
\end{chapterexercises}

\chapter{Sliding-window attention and cache policies}
\label{ch:sliding-window}
Sliding-window attention limits visible history and changes cache semantics. \citep{jiang2023mistral}

\[
\operatorname{visible}(t,W)=\{\max(0,t+1-W),\ldots,t\}.
\]

\begin{intuitionbox}{Intuition}
A cache policy is model semantics.
\end{intuitionbox}

\begin{definitionbox}{Mathematical or contract statement}
For query position \(t\), windowed attention exposes only the bounded suffix \(\operatorname{visible}(t,W)\). Append-only attention is the special case where \(W\) is at least the full accumulated sequence length.
\end{definitionbox}

\begin{implementationbox}{Implementation interpretation}
Keep phase-0 append-only; add sliding windows as a later policy.
\end{implementationbox}

\begin{verificationbox}{How to verify this works}
\begin{itemize}
\item Check shapes and dtypes at the boundary.
\item Compare deterministic outputs against PyTorch golden vectors.
\item Record tolerance, seed, model artifact, and backend.
\item Do not benchmark until validation passes.
\end{itemize}
\end{verificationbox}

\begin{warningbox}{Common failure modes}
\begin{itemize}
\item Letting framework defaults choose dtype or layout.
\item Testing only a batch size of one.
\item Depending on upstream checkpoint names in runtime code.
\item Treating benchmark output as validation.
\end{itemize}
\end{warningbox}

\begin{chapterexercises}
\item State the central invariant in one sentence.\\
\textit{Hint.} Use the conceptual guide vocabulary.
\item Construct a toy example with small shapes.\\
\textit{Hint.} Use $B=2$, $T=3$, $H=8$ where possible.
\item Name a validation test that would catch a plausible bug.\\
\textit{Hint.} Prefer generated golden vectors to handwritten long arrays.
\end{chapterexercises}

\chapter{Modern dense model-family adapters}
\label{ch:family-adapters}
Modern dense families differ in norm, attention, tokenizer, position, and checkpoint naming details.



\begin{intuitionbox}{Intuition}
Adapters translate a family into internal contracts.
\end{intuitionbox}

\begin{definitionbox}{Mathematical or contract statement}
A family adapter supplies config rules, tensor mappings, and validation plans.
\end{definitionbox}

\begin{implementationbox}{Implementation interpretation}
Keep adapters thin so core runtime remains reusable.
\end{implementationbox}

\begin{verificationbox}{How to verify this works}
\begin{itemize}
\item Check shapes and dtypes at the boundary.
\item Compare deterministic outputs against PyTorch golden vectors.
\item Record tolerance, seed, model artifact, and backend.
\item Do not benchmark until validation passes.
\end{itemize}
\end{verificationbox}

\begin{warningbox}{Common failure modes}
\begin{itemize}
\item Letting framework defaults choose dtype or layout.
\item Testing only a batch size of one.
\item Depending on upstream checkpoint names in runtime code.
\item Treating benchmark output as validation.
\end{itemize}
\end{warningbox}

\begin{chapterexercises}
\item State the central invariant in one sentence.\\
\textit{Hint.} Use the conceptual guide vocabulary.
\item Construct a toy example with small shapes.\\
\textit{Hint.} Use $B=2$, $T=3$, $H=8$ where possible.
\item Name a validation test that would catch a plausible bug.\\
\textit{Hint.} Prefer generated golden vectors to handwritten long arrays.
\end{chapterexercises}

\chapter{Qwen/Gemma/Mistral-style implementation considerations}
\label{ch:qwen-gemma-mistral}
Qwen, Gemma, and Mistral-style models illustrate why family adapters are necessary. \citep{qwen2024qwen25,gemma2024,jiang2023mistral}



\begin{intuitionbox}{Intuition}
Model-family names are not architecture contracts.
\end{intuitionbox}

\begin{definitionbox}{Mathematical or contract statement}
The common core stays decoder inference; adapters specify variant details.
\end{definitionbox}

\begin{implementationbox}{Implementation interpretation}
Do not claim support without converter tests and golden vectors.
\end{implementationbox}

\begin{verificationbox}{How to verify this works}
\begin{itemize}
\item Check shapes and dtypes at the boundary.
\item Compare deterministic outputs against PyTorch golden vectors.
\item Record tolerance, seed, model artifact, and backend.
\item Do not benchmark until validation passes.
\end{itemize}
\end{verificationbox}

\begin{warningbox}{Common failure modes}
\begin{itemize}
\item Letting framework defaults choose dtype or layout.
\item Testing only a batch size of one.
\item Depending on upstream checkpoint names in runtime code.
\item Treating benchmark output as validation.
\end{itemize}
\end{warningbox}

\begin{chapterexercises}
\item State the central invariant in one sentence.\\
\textit{Hint.} Use the conceptual guide vocabulary.
\item Construct a toy example with small shapes.\\
\textit{Hint.} Use $B=2$, $T=3$, $H=8$ where possible.
\item Name a validation test that would catch a plausible bug.\\
\textit{Hint.} Prefer generated golden vectors to handwritten long arrays.
\end{chapterexercises}

\chapter{Long-context validation}
\label{ch:long-context}
Long-context validation stresses position offsets, cache memory, visibility policies, and tolerance accumulation.



\begin{intuitionbox}{Intuition}
Context length is a semantic claim and a systems claim.
\end{intuitionbox}

\begin{definitionbox}{Mathematical or contract statement}
Long-context tests cover positions near limits, cache lengths, memory estimates, and decode equivalence.
\end{definitionbox}

\begin{implementationbox}{Implementation interpretation}
Use generated vectors and summaries rather than hard-coded long arrays.
\end{implementationbox}

\begin{verificationbox}{How to verify this works}
\begin{itemize}
\item Check shapes and dtypes at the boundary.
\item Compare deterministic outputs against PyTorch golden vectors.
\item Record tolerance, seed, model artifact, and backend.
\item Do not benchmark until validation passes.
\end{itemize}
\end{verificationbox}

\begin{warningbox}{Common failure modes}
\begin{itemize}
\item Letting framework defaults choose dtype or layout.
\item Testing only a batch size of one.
\item Depending on upstream checkpoint names in runtime code.
\item Treating benchmark output as validation.
\end{itemize}
\end{warningbox}

\begin{chapterexercises}
\item State the central invariant in one sentence.\\
\textit{Hint.} Use the conceptual guide vocabulary.
\item Construct a toy example with small shapes.\\
\textit{Hint.} Use $B=2$, $T=3$, $H=8$ where possible.
\item Name a validation test that would catch a plausible bug.\\
\textit{Hint.} Prefer generated golden vectors to handwritten long arrays.
\end{chapterexercises}

\chapter{Quantisation overview, marked as later work}
\label{ch:quantisation}
Quantisation reduces memory bandwidth and storage but is later work relative to phase-0 float32 correctness.



\begin{intuitionbox}{Intuition}
Quantisation changes the numerical contract.
\end{intuitionbox}

\begin{definitionbox}{Mathematical or contract statement}
Quantised inference requires scale or block metadata and a documented comparison tolerance.
\end{definitionbox}

\begin{implementationbox}{Implementation interpretation}
Do not let quantisation leak into phase-0 artifacts.
\end{implementationbox}

\begin{verificationbox}{How to verify this works}
\begin{itemize}
\item Check shapes and dtypes at the boundary.
\item Compare deterministic outputs against PyTorch golden vectors.
\item Record tolerance, seed, model artifact, and backend.
\item Do not benchmark until validation passes.
\end{itemize}
\end{verificationbox}

\begin{warningbox}{Common failure modes}
\begin{itemize}
\item Letting framework defaults choose dtype or layout.
\item Testing only a batch size of one.
\item Depending on upstream checkpoint names in runtime code.
\item Treating benchmark output as validation.
\end{itemize}
\end{warningbox}

\begin{chapterexercises}
\item State the central invariant in one sentence.\\
\textit{Hint.} Use the conceptual guide vocabulary.
\item Construct a toy example with small shapes.\\
\textit{Hint.} Use $B=2$, $T=3$, $H=8$ where possible.
\item Name a validation test that would catch a plausible bug.\\
\textit{Hint.} Prefer generated golden vectors to handwritten long arrays.
\end{chapterexercises}
